\section{Forelesning 6: \emph{Digitale økosystemer og plattformer}}
\label{sec:lec06}
\settermlabel{terms:lec06}

\subsection*{Oversikt}
Forelesning 6 (Sebastian Winslow) tar for seg \emph{digitale økosystemer} som
en strategisk organiseringsform som skiller seg fra den klassiske verdikjeden.
Hovedpoenget er at digital teknologi muliggjør \emph{nettverksbasert}
verdiskaping på tvers av uavhengige aktører, styrt av plattformregler heller
enn eierskap. Dette knytter direkte til Weill \& Woerners typer
(\autoref{sec:lec03}) -- særlig \emph{ecosystem driver} og \emph{modular
producer}.

\subsection{Hovedteori}

\subsubsection*{Hva er et digitalt økosystem?}

\begin{definition}{Digitalt økosystem}
Et nettverk av gjensidig avhengige aktører (virksomheter, kunder, partnere,
utviklere) som samskaper og fanger verdi sammen, ofte formidlet av en
\emph{plattform}. Kjennetegn: nettverkseffekter, plattformstyring,
samskapelse av verdi og dataflyt på tvers av aktørene.
\end{definition}

\subsubsection*{Verdikjede vs.\ økosystem}
\begin{itemize}
  \item \term{Value chain} -- lineær, kontrollert, sekvensiell produksjon.
    Virksomheten eier eller koordinerer innsatsfaktorer $\rightarrow$
    transformasjon $\rightarrow$ utdata.
  \item \term{Ecosystem} -- nettverksbasert, flere uavhengige aktører,
    ikke-lineær verdiskaping, styrt av plattformregler heller enn eierskap.
\end{itemize}

\keypoint{Den strategiske implikasjonen: i et økosystem flyttes fortrinnet
fra \emph{å eie aktivitetene} til \emph{å styre plattformen} som koordinerer dem.}

\subsubsection*{Plattformer som økosystemhub}
Plattformer er den vanligste formen for økosystemhub. Sentrale egenskaper:
\begin{itemize}
  \item \term{Multi-sided platform} -- betjener to eller flere kundegrupper
    samtidig (f.eks.\ kunder og selgere på Amazon).
  \item \modelref{network-effects}{Network effects} -- verdien for én bruker vokser med antall
    andre brukere.
  \item \term{Platform envelopment} -- en plattform utvider sitt domene ved å
    absorbere tilstøtende plattformer.
  \item \term{Tipping} -- markeder med sterke nettverkseffekter har en
    tendens til å \guillemotleft{}tippe\guillemotright{} mot én dominerende plattform.
\end{itemize}

\subsubsection*{Ubers \guillemotleft{}napkin\guillemotright{}: virtuous cycle}
Den hånd-tegnede Uber-flywheelen viser logikken bak nettverkseffekter:
\emph{flere kunder $\rightarrow$ flere sjåfører $\rightarrow$ større
geografisk dekning $\rightarrow$ raskere pickup / mindre nedetid for
sjåfører $\rightarrow$ lavere priser $\rightarrow$ flere kunder.}

\begin{figure}[H]
\centering
\includegraphics[width=0.7\linewidth]{lec06-uber-network-effects-napkin.png}
\caption{\guillemotleft{}The Napkin\guillemotright{} -- Ubers virtuous-cycle av nettverkseffekter.}
\end{figure}

\subsubsection*{\modelref{network-effects}{Aksene av nettverkseffekter}}
Nettverkseffekter er ikke én ting; de er fire klasser. Aksene:
\emph{Same-side vs.\ Cross-side} $\times$ \emph{Positive vs.\ Negative}. En
plattform må forstå hvilke effekter som faktisk er i spill -- og hvilke som
er negative (overbelastning, kongestion, ulike interesser mellom sider).

\begin{figure}[H]
\centering
\includegraphics[width=0.8\linewidth]{lec06-axes-of-network-effects.png}
\caption{De fire klassene av nettverkseffekter: Same-side / Cross-side
$\times$ Positive / Negative.}
\end{figure}

\subsubsection*{\guillemotleft{}Three-dimensional chess\guillemotright{}: plattformkonkurranse}
Parker et al.: plattformer spiller samtidig tre lag av konkurranse --
\emph{plattform vs.\ plattform}, \emph{plattform vs.\ partner}, og
\emph{partner vs.\ partner}. Plattformer vokser pien heller enn å re-fordele
den; det er derfor tradisjonell Five Forces-tenkning ofte undervurderer
plattformmakt.

\begin{figure}[H]
\centering
\includegraphics[width=0.85\linewidth]{lec06-three-dimensional-chess-platform-competition.png}
\caption{\guillemotleft{}Three-dimensional chess\guillemotright{} -- tre samtidige nivåer av
plattformkonkurranse.}
\end{figure}

\subsubsection*{Styring av økosystemer}
\begin{itemize}
  \item \term{Platform governance} -- regler for hvem som får delta, hvilken
    data som deles, hvordan verdi fordeles.
  \item \term{Openness} vs.\ \term{closedness} -- avveining mellom å åpne for
    flere tredjeparter (mer verdi) og å beholde kontroll (mer fangst).
  \item \term{Boundary resources} -- API-er, SDK-er, dokumentasjon som lar
    tredjeparter koble seg på.
\end{itemize}

\subsubsection*{Verdiskaping og verdifangst i økosystem}
\begin{itemize}
  \item \term{Value co-creation} -- verdi oppstår gjennom samspillet, ikke i én
    aktørs aktiviteter.
  \item \term{Value capture} -- plattformen tar typisk en andel (provisjon,
    abonnement, datatilgang).
  \item Sentral spenning: hvor mye verdi skal plattformen fange før
    deltakerne forlater den?
\end{itemize}

\subsubsection*{Platform Business Model Canvas (PBMC)}
PBMC tilpasser Osterwalder-canvas til plattformlogikken: separate
Producer- og Consumer-segmenter, eksplisitte felt for \emph{Interaction},
\emph{Stimuli}, \emph{Facilitation}, og \emph{Monetization} -- og for
\emph{governance}.

\begin{figure}[H]
\centering
\includegraphics[width=0.95\linewidth]{lec06-platform-business-model-canvas-compact.png}
\caption{Platform Business Model Canvas -- kompakt versjon.}
\end{figure}

\begin{figure}[H]
\centering
\includegraphics[width=0.95\linewidth]{lec06-platform-business-model-canvas-details.png}
\caption{PBMC -- detaljvisning: (1) Platform Users / Stakeholders, (2) Value
Proposition (Job / Pain / Gain / Core Value Unit / Transaction), (3) Access
and Usage (Resources, Activities, Channels, Filter, Promotion, Governance).}
\end{figure}

\subsection{Sentrale fagbegreper}\label{terms:lec06}
\begin{itemize}
  \item \term{Digital ecosystem}, \term{platform}, \term{network effects}.
  \item \term{Multi-sided platform}, \term{boundary resources},
    \term{platform envelopment}, \term{tipping}.
  \item \term{Same-side} vs.\ \term{cross-side} effects (positive/negative).
  \item \term{Platform governance}, \term{openness/closedness}.
  \item \term{Value co-creation}, \term{value capture}.
  \item \term{Complementors} -- tredjeparter som bygger på plattformen.
  \item \term{Three-dimensional chess} -- de tre konkurransenivåene.
\end{itemize}

\subsection{Diagnostisk rammeverk for økosystemstrategi}
\begin{enumerate}
  \item Hvor er nettverkseffektene -- direkte, indirekte, ingen?
    Same-/cross-side? Positive/negative?
  \item Hvem er hubben -- vi, en konkurrent, en plattformeier?
  \item Hva er boundary resources -- åpne nok til å tiltrekke
    komplementorer, lukkede nok til å beholde fangst?
  \item Hvordan fordeles verdi -- og er fordelingen stabil over tid?
  \item Hvilke nivåer av \guillemotleft{}three-dimensional chess\guillemotright{} spilles -- og hvor er vi?
\end{enumerate}

\subsection{Eksempler og caser}
\begin{eksempelboks}{Amazon, Alibaba}
\emph{Ecosystem drivers} med svært bred kundekunnskap og bred koordinering
av tredjeparter. Marketplace + AWS + Prime + logistikk = pien vokser.
\end{eksempelboks}

\begin{eksempelboks}{Apple iOS}
Plattform med streng styring av komplementorer (App Store-regler,
30~\% provisjon). Klassisk eksempel på \emph{closedness} med høy verdifangst,
men også på regulatorisk press når makten blir for stor.
\end{eksempelboks}

\begin{eksempelboks}{Lobyco som modular producer}
Potensiell \emph{modular producer} som plugger inn
lojalitets-/betalingsinfrastruktur i \emph{andres} økosystemer. Dette er et
begrenset økosystem-spill -- Coop blir leverandør, ikke hub.
\end{eksempelboks}

\begin{eksempelboks}{Coop som ecosystem driver?}
Ville kreve å være primær hub for danske forbrukeres matplanlegging og
hverdagslogistikk. Ingen evidens støtter denne ambisjonen i dag; lukkingen
av Coop.dk og fokuset på butikk-først peker motsatt vei.
\end{eksempelboks}

\subsection{Eksamensrelevans}
\begin{itemize}
  \item Definere økosystem og kontrastere mot verdikjede.
  \item Forklare ulike typer nettverkseffekter med eksempler -- ikke bare
    \guillemotleft{}det er nettverkseffekter\guillemotright{}.
  \item Diskutere platform governance: åpenhet vs.\ kontroll, og avveiningen
    mellom verdiskaping og verdifangst.
  \item Bruke PBMC når en konkret plattform analyseres.
  \item Plassere en konkret virksomhet (Coop, Lobyco) i økosystem-/plattform-
    terminologi -- og være presis om hva som er evidens-basert.
  \item Knytte til \autoref{sec:lec03} (DBM-typer): \emph{ecosystem driver} og
    \emph{modular producer} bor i økosystemverdenen.
\end{itemize}
